% =============================================================================
%  Definiciones (\newdefinition de la clase) y macros de redacción
% =============================================================================
\usepackage{xcolor}
% Marca visible para lo que falta redactar o decidir. Eliminar todas antes de entregar:
%   grep -rn "pendiente" chapters appendices front-matter
\newcommand{\pendiente}[1]{\textcolor{red}{\textbf{[PENDIENTE:} #1\textbf{]}}}

\newdefinition{backend}{backend}{backends}%
  {Parte del sistema responsable de la lógica de negocio, la persistencia y la
   exposición de servicios, sin contacto directo con la interfaz de usuario}
\newdefinition{frontend}{frontend}{frontends}%
  {Parte del sistema responsable de la interacción con el usuario y de la
   presentación de la información}
\newdefinition{instancia}{instancia de examen}{instancias de examen}%
  {Examen escaneado de un estudiante concreto, con sus páginas, anotaciones y notas}
\newdefinition{rubrica}{rúbrica}{rúbricas}%
  {Conjunto de criterios con puntuación con los que se califica un problema}
\newdefinition{mock}{simulación del backend}{simulaciones del backend}%
  {Implementación en el navegador que responde a las peticiones con los mismos
   formatos que el backend real, usada para desarrollar y probar sin él}
\newdefinition{agente}{agente de programación}{agentes de programación}%
  {Sistema basado en un modelo de lenguaje que lee y modifica código, ejecuta
   órdenes y verifica resultados con un grado de autonomía configurable}
\newdefinition{worktree}{worktree}{worktrees}%
  {Directorio de trabajo adicional de Git asociado a otra rama del mismo
   repositorio, que permite trabajar en paralelo sin compartir ficheros}

% Fragmentos de código (JavaScript/JSX) en el capítulo de implementación
\usepackage{listings}
\lstdefinelanguage{JavaScript}{%
  keywords={const,let,var,function,return,if,else,for,of,new,async,await,try,catch,
            finally,throw,export,import,from,default,null,true,false,typeof},
  sensitive=true,
  comment=[l]{//},
  morecomment=[s]{/*}{*/},
  morestring=[b]',
  morestring=[b]",
  morestring=[b]`}
\lstset{language=JavaScript, basicstyle=\ttfamily\footnotesize,
        keywordstyle=\bfseries, commentstyle=\itshape\color{gray},
        columns=fullflexible, keepspaces=true, breaklines=true,
        frame=single, framerule=0.3pt, rulecolor=\color{gray},
        captionpos=b, abovecaptionskip=4pt, xleftmargin=2pt, xrightmargin=2pt}
\renewcommand{\lstlistingname}{Código}
\renewcommand{\lstlistlistingname}{Índice de códigos}
